% ==============================================================
\section{Application: Discounted Utility}
\label{sec:discounted-utility}
% ==============================================================

We now apply the recursive structure to intertemporal choice. The standard discounted utility model---the workhorse of dynamic economics---turns out to be a monotone transformation of a generalized mean. This connection explains why dynamic problems with CES preferences are tractable.

% --------------------------------------------------------------
\subsection{The Standard Formulation}
\label{subsec:standard-formulation}
% --------------------------------------------------------------

Consider an agent evaluating a consumption stream $\{c_t\}_{t=0}^{T}$ over a finite horizon. The standard objective is:
\begin{equation}
V = \sum_{t=0}^{T} \beta^t u(c_t),
\label{eq:standard-objective}
\end{equation}
where $\beta \in (0,1)$ is the discount factor and $u(\cdot)$ is the period utility function.

\paragraph{Power Utility.}
The most common specification is \emph{power utility} (also called isoelastic or CRRA utility):
\begin{equation}
u(c) = \frac{c^{1-1/\theta} - 1}{1 - 1/\theta},
\label{eq:power-utility}
\end{equation}
where $\theta > 0$ is the \emph{intertemporal elasticity of substitution} (IES). This measures the willingness to substitute consumption across time in response to changes in the interest rate.

\begin{calloutnote}[The Role of $-1$]
The term $-1$ in the numerator ensures continuity at $\theta = 1$. As $\theta \to 1$, L'H\^opital's rule gives:
\begin{equation}
\lim_{\theta \to 1} \frac{c^{1-1/\theta} - 1}{1 - 1/\theta} = \log c.
\end{equation}
Without the $-1$, the limit would not exist. This is a normalization that affects levels but not optimal choices.
\end{calloutnote}

The full objective becomes:
\begin{equation}
V = \sum_{t=0}^{T} \beta^t \frac{c_t^{1-1/\theta} - 1}{1 - 1/\theta}.
\label{eq:full-objective}
\end{equation}

% --------------------------------------------------------------
\subsection{Transformation to a Generalized Mean}
\label{subsec:transformation-mean}
% --------------------------------------------------------------

We now show that $V$ is a monotone transformation of a classical generalized mean. The key steps are: (i) normalize the weights, (ii) separate the constant terms, and (iii) apply a monotone transformation.

\paragraph{Step 1: Normalizing the Weights.}
The discount factors $\{\beta^t\}_{t=0}^{T}$ do not sum to one. Define the normalization constant:
\begin{equation}
\kappa^{-1} = \sum_{t=0}^{T} \beta^t = \frac{1 - \beta^{T+1}}{1 - \beta}.
\label{eq:kappa-def}
\end{equation}
Then $\bar{\omega}_t \equiv \kappa \beta^t$ satisfies $\sum_{t=0}^{T} \bar{\omega}_t = 1$, forming valid weights for a generalized mean.

\paragraph{Step 2: Separating the Constant.}
Multiply and divide the objective by $\kappa$:
\begin{equation}
V = \kappa^{-1} \sum_{t=0}^{T} \bar{\omega}_t \frac{c_t^{1-1/\theta} - 1}{1 - 1/\theta} = \kappa^{-1} \left[ \frac{1}{1 - 1/\theta} \sum_{t=0}^{T} \bar{\omega}_t c_t^{1-1/\theta} - \frac{1}{1 - 1/\theta} \right].
\label{eq:separate-constant}
\end{equation}
The second term uses $\sum_t \bar{\omega}_t = 1$. Rearranging:
\begin{equation}
\kappa V = \frac{1}{1 - 1/\theta} \left[ \sum_{t=0}^{T} \bar{\omega}_t c_t^{1-1/\theta} - 1 \right].
\label{eq:rearranged}
\end{equation}

\paragraph{Step 3: Isolating the Inner Sum.}
Define $\rho \equiv 1 - 1/\theta$, the power parameter. Then:
\begin{equation}
\kappa (1 - 1/\theta) V + 1 = \sum_{t=0}^{T} \bar{\omega}_t c_t^{\rho} = \bar{\Ssum}(\bc; \bar{\bom}, \rho),
\label{eq:inner-sum-isolated}
\end{equation}
where $\bar{\Ssum}$ is the inner sum of the classical generalized mean and $\bc = (c_0, \ldots, c_T)$.

\paragraph{Step 4: Applying a Monotone Transformation.}
The inner sum relates to the classical mean by $\bar{\Ssum} = \Mbar^\rho$. Taking the $1/\rho$ power of both sides:
\begin{equation}
\left( \kappa (1 - 1/\theta) V + 1 \right)^{1/\rho} = \Mbar(\bc; \bar{\bom}, \rho).
\label{eq:monotone-transform}
\end{equation}

\begin{proposition}[Discounted Utility as a Generalized Mean]
\label{prop:utility-as-mean}
The standard discounted power utility $V = \sum_{t=0}^{T} \beta^t u(c_t)$ is a monotone transformation of the classical generalized mean:
\begin{equation}
V = \frac{1}{\kappa(1-1/\theta)} \left[ \Mbar(\bc; \bar{\bom}, \rho)^{\rho} - 1 \right],
\end{equation}
where $\rho = 1 - 1/\theta$, the weights are $\bar{\omega}_t = \kappa \beta^t$ with $\kappa = (1-\beta)/(1-\beta^{T+1})$, and $\Mbar(\bc; \bar{\bom}, \rho) = \left( \sum_{t=0}^{T} \bar{\omega}_t c_t^\rho \right)^{1/\rho}$.
\end{proposition}

\begin{calloutimportant}[Why This Matters]
Since utility is ordinal, maximizing $V$ is equivalent to maximizing $\Mbar(\bc; \bar{\bom}, \rho)$. All the machinery developed for generalized means---recursive decomposition, sufficient statistics, backward induction---applies directly to intertemporal choice.
\end{calloutimportant}

% --------------------------------------------------------------
\subsection{The Canonical Form}
\label{subsec:canonical-intertemporal}
% --------------------------------------------------------------

We can also express discounted utility in canonical form. Recall from the transformation theory that:
\begin{equation}
\Mbar(\bc; \bar{\bom}, \rho) = \lambda(\bar{\bom}, \rho)^{-1} \M(\bc; \bom, \rho),
\label{eq:classical-canonical-relation}
\end{equation}
where $\bom$ are the canonical weights derived from $\bar{\bom}$ and $\lambda$ is the scaling factor.

\paragraph{Deriving the Canonical Weights.}
The canonical weights satisfy:
\begin{equation}
\omega_t = \frac{\bar{\omega}_t^{1/(1-\rho)}}{\sum_{s=0}^{T} \bar{\omega}_s^{1/(1-\rho)}} = \frac{(\kappa \beta^t)^{\theta}}{\sum_{s=0}^{T} (\kappa \beta^s)^{\theta}} = \frac{\beta^{\theta t}}{\sum_{s=0}^{T} \beta^{\theta s}}.
\label{eq:canonical-weights}
\end{equation}
Define $\tilde{\beta} \equiv \beta^\theta$. Then:
\begin{equation}
\omega_t = \frac{\tilde{\beta}^t}{\sum_{s=0}^{T} \tilde{\beta}^s} = \tilde{\kappa} \tilde{\beta}^t,
\label{eq:canonical-weights-simple}
\end{equation}
where $\tilde{\kappa} = (1 - \tilde{\beta})/(1 - \tilde{\beta}^{T+1})$.

\begin{calloutnote}[Effective Discount Factor]
The canonical weights use an \emph{effective discount factor} $\tilde{\beta} = \beta^\theta$ that combines the primitive discount factor $\beta$ with the intertemporal elasticity $\theta$. When $\theta = 1$ (log utility), $\tilde{\beta} = \beta$. When $\theta > 1$ (high substitutability), $\tilde{\beta} < \beta$---the agent behaves as if more impatient. When $\theta < 1$, $\tilde{\beta} > \beta$.
\end{calloutnote}

\paragraph{The Canonical Representation.}
In canonical form, discounted utility becomes:
\begin{equation}
\M(\bc; \bom, \rho) = \left( \sum_{t=0}^{T} \omega_t^{1-\rho} c_t^\rho \right)^{1/\rho} = \left( \sum_{t=0}^{T} (\tilde{\kappa} \tilde{\beta}^t)^{1-\rho} c_t^\rho \right)^{1/\rho}.
\label{eq:canonical-utility}
\end{equation}
This form will prove useful when we analyze price indices and expenditure shares in intertemporal settings.

% --------------------------------------------------------------
\subsection{Recursive Formulation}
\label{subsec:recursive-utility}
% --------------------------------------------------------------

The recursive decomposition from \cref{sec:recursive-decomposition} applies directly. Since discounted utility is a generalized mean, it inherits the recursive structure automatically.

\paragraph{Applying the Decomposition.}
Let $\V_t$ denote the utility from date $t$ onward, expressed as a generalized mean of $\{c_t, c_{t+1}, \ldots, c_T\}$. From \cref{prop:recursive-decomposition}:
\begin{equation}
\V_t = \Mbar\Big( c_t, \V_{t+1}; \bar{\omega}_t^{(t)}, \rho \Big),
\label{eq:bellman-mean}
\end{equation}
where $\bar{\omega}_t^{(t)}$ is the renormalized weight on current consumption and $\V_{t+1}$ is the continuation value.

\paragraph{Computing the Renormalized Weight.}
In the original problem, the weight on $c_t$ is $\bar{\omega}_t = \kappa \beta^t$, and the sum of weights from $t$ onward is:
\begin{equation}
\sum_{s=t}^{T} \bar{\omega}_s = \kappa \beta^t \frac{1 - \beta^{T+1-t}}{1 - \beta}.
\label{eq:weight-sum-from-t}
\end{equation}
The renormalized weight on $c_t$ is:
\begin{equation}
\bar{\omega}_t^{(t)} = \frac{\bar{\omega}_t}{\sum_{s=t}^{T} \bar{\omega}_s} = \frac{1 - \beta}{1 - \beta^{T+1-t}}.
\label{eq:renormalized-current}
\end{equation}

\paragraph{The Infinite-Horizon Limit.}
As $T \to \infty$, the renormalized weight converges to:
\begin{equation}
\bar{\omega}_t^{(t)} \to 1 - \beta.
\label{eq:weight-infinite}
\end{equation}
This is independent of $t$---the problem becomes stationary. The recursion simplifies to:
\begin{equation}
\V_t = \Mbar\Big( c_t, \V_{t+1}; 1-\beta, \rho \Big) = \left( (1-\beta) c_t^\rho + \beta \V_{t+1}^\rho \right)^{1/\rho}.
\label{eq:bellman-infinite}
\end{equation}

\begin{proposition}[Recursive Discounted Utility]
\label{prop:recursive-utility}
In the infinite-horizon limit, discounted power utility satisfies the recursion:
\begin{equation}
\V_t = \left( (1-\beta) c_t^\rho + \beta \V_{t+1}^\rho \right)^{1/\rho},
\end{equation}
where $\rho = 1 - 1/\theta$. The continuation value $\V_{t+1}$ is a sufficient statistic for all future consumption.
\end{proposition}

\paragraph{Terminal Condition.}
For a finite horizon, we need a terminal condition. At $t = T$:
\begin{equation}
\V_T = c_T.
\label{eq:terminal}
\end{equation}
We then solve backward: $\V_{T-1} = \Mbar(c_{T-1}, \V_T; \bar{\omega}_{T-1}^{(T-1)}, \rho)$, and so on.

% --------------------------------------------------------------
\subsection{What Does Not Work}
\label{subsec:what-does-not-work}
% --------------------------------------------------------------

The recursive structure relies on the geometric discount weights $\beta^t$. Not all weight structures admit this decomposition.

\paragraph{The Quasi-Hyperbolic Case.}
Consider the $\beta$-$\delta$ preferences common in behavioral economics:
\begin{equation}
V = u(c_0) + \beta \sum_{t=1}^{T} \delta^t u(c_t),
\label{eq:beta-delta}
\end{equation}
where $\beta \in (0,1)$ captures "present bias" and $\delta \in (0,1)$ is the long-run discount factor. The weights are:
\begin{equation}
\text{weight on } c_0 = 1, \qquad \text{weight on } c_t = \beta \delta^t \text{ for } t \geq 1.
\label{eq:beta-delta-weights}
\end{equation}

\paragraph{Why the Recursion Fails.}
Attempt to write this recursively by separating $c_0$:
\begin{equation}
V = u(c_0) + \beta \delta \sum_{t=1}^{T} \delta^{t-1} u(c_t) = u(c_0) + \beta \delta \cdot V',
\label{eq:beta-delta-attempt}
\end{equation}
where $V' = \sum_{t=1}^{T} \delta^{t-1} u(c_t)$. But from the perspective of period 1, the agent evaluates:
\begin{equation}
V_1 = u(c_1) + \beta \sum_{t=2}^{T} \delta^{t-1} u(c_t) \neq V'.
\label{eq:beta-delta-v1}
\end{equation}
The object $V'$ (which discounts geometrically from period 1) is \emph{not the same} as $V_1$ (which applies the present-bias factor $\beta$ again at period 1). The inner object differs from the outer object.

\begin{calloutnote}[The Problem]
With $\beta$-$\delta$ preferences, the continuation value $V'$ that appears in today's recursion is not the same as the value function $V_1$ that the agent will actually use tomorrow. The weights cannot be factored to produce a self-similar structure.

This is why $\beta$-$\delta$ preferences lead to time-inconsistent behavior: the agent's future self solves a different problem than the current self anticipates. We do not pursue this here, but it illustrates that the recursive structure of standard discounted utility is a substantive restriction, not merely a notational convenience.
\end{calloutnote}

\paragraph{The General Lesson.}
For the recursion to work, the weight structure must satisfy:
\begin{equation}
\frac{\omega_t}{\sum_{s=t}^{T} \omega_s} = f(\text{parameters only}),
\label{eq:factorization-condition}
\end{equation}
independent of $t$ (at least in the infinite-horizon limit). Geometric weights $\omega_t \propto \beta^t$ satisfy this with $f = 1 - \beta$. Quasi-hyperbolic weights do not.

% --------------------------------------------------------------
\subsection{Summary}
\label{subsec:discounted-summary}
% --------------------------------------------------------------

\begin{calloutimportant}[Key Results]
\begin{enumerate}[label=(\roman*), leftmargin=2em]
    \item Standard discounted power utility is a \textbf{monotone transformation} of a classical generalized mean with weights $\bar{\omega}_t = \kappa \beta^t$.
    
    \item The \textbf{canonical form} uses an effective discount factor $\tilde{\beta} = \beta^\theta$, combining impatience and substitutability.
    
    \item The \textbf{recursive formulation} is $\V_t = \Mbar(c_t, \V_{t+1}; 1-\beta, \rho)$ in the infinite-horizon limit, where $\rho = 1 - 1/\theta$.
    
    \item The \textbf{continuation value} $\V_{t+1}$ is a sufficient statistic for all future consumption---this is why dynamic programming works with CES preferences.
\end{enumerate}
\end{calloutimportant}

\begin{callouttip}[Exercise]
\begin{enumerate}[label=(\alph*), leftmargin=1.5em]
    \item Verify that when $\theta = 1$ (log utility), the recursion $\V_t = ((1-\beta)c_t^\rho + \beta \V_{t+1}^\rho)^{1/\rho}$ reduces to $\V_t = c_t^{1-\beta} \V_{t+1}^\beta$.
    \item For the finite-horizon case with $T = 2$, compute $\V_0$ explicitly as a function of $(c_0, c_1, c_2)$ using backward induction.
    \item Show that the result in (b) equals the classical mean $\Mbar((c_0, c_1, c_2); (\bar{\omega}_0, \bar{\omega}_1, \bar{\omega}_2), \rho)$.
\end{enumerate}
\end{callouttip}
